% Lean: EconLean/Mechanism/{Revelation,VCG,Myerson,OptimalMechanism,BilateralTrade}.lean
\chapter{Mechanism Design}
\label{chap:mechanism}

Source: MWG Ch.~23; Myerson (1981).

\section{Implementation and the revelation principle}
\label{sec:revelation}

\begin{definition}[Dominant-strategy equilibrium, DSIC, BIC; MWG Defs.~23.B.5, 23.B.6, 23.D.1]
\label{def:dsic}
\lean{Econ.IsDominantEq, Econ.IsDSIC, Econ.IsBIC, Econ.mechGame}
\leanok
\uses{def:bne}
A mechanism $g : \prod_i M_i \to X$ with strategies $s_i : \Theta_i \to M_i$ is in dominant-strategy equilibrium if
$s_i(\theta_i)$ is optimal for type $\theta_i$ against every message profile of the others. A social choice function
$f$ is DSIC if truth-telling is optimal against every report of the others, and BIC for prior $p$ if for every
$i, \theta_i, \theta_i'$: $\sum_{\theta : \theta_i \text{ fixed}} p(\theta)\, u_i(f(\theta_i', \theta_{-i}), \theta_i) \le
\sum_{\theta : \theta_i \text{ fixed}} p(\theta)\, u_i(f(\theta), \theta_i)$.
\end{definition}

\begin{theorem}[Revelation principle; MWG Props.~23.C.1, 23.D.1]
\label{thm:revelation}
\lean{Econ.play_update', Econ.revelation_dominant, Econ.revelation_bayesian}
\leanok
\uses{def:dsic, thm:bne}
If $s$ is a dominant-strategy equilibrium (resp.\ Bayesian Nash equilibrium) of $g$, then
$f = g \circ s$ is DSIC (resp.\ BIC).
\end{theorem}
\begin{proof}
\leanok
A report $\theta_i'$ in the direct mechanism has the same effect as the message $s_i(\theta_i')$ in $g$, which is not
profitable. For BNE use the interim characterisation.
\end{proof}

\section{Groves and VCG mechanisms}
\label{sec:vcg}

\begin{definition}[Groves and Clarke transfers]
\label{def:groves}
\lean{Econ.IsEfficientRule, Econ.IgnoresOwn, Econ.grovesTransfer, Econ.othersMax, Econ.clarkeTransfer, Econ.auctionValue, Econ.highBidder}
\leanok
$k$ is efficient if it maximizes $\sum_i v_i(x,\theta_i)$. The Groves transfer is
$t_i(\theta) = \sum_{j\ne i} v_j(k(\theta),\theta_j) + h_i(\theta_{-i})$; Clarke takes
$h_i = -\max_x \sum_{j\ne i} v_j(x,\theta_j)$.
\end{definition}

\begin{theorem}[Groves mechanisms are strategy-proof; MWG Props.~23.C.4--23.C.5]
\label{thm:groves}
\lean{Econ.groves_truthful, Econ.othersMax_ignoresOwn, Econ.clarke_truthful, Econ.clarke_transfer_nonpos, Econ.clarke_individuallyRational}
\leanok
\uses{def:groves}
Truth-telling is dominant in every Groves mechanism; Clarke transfers are nonpositive, and with
nonnegative valuations truthful agents get nonnegative payoffs.
\end{theorem}
\begin{proof}
\leanok
Agent $i$'s payoff is the reported total value of $k(\theta_i', \theta_{-i})$ evaluated at $i$'s true type,
plus a term independent of $\theta_i'$; efficiency at the true report maximizes it.
\end{proof}

\section{Myerson's lemma and revenue equivalence}
\label{sec:myerson}

\begin{lemma}[Integral from a monotone sandwich]
\label{lem:sandwich}
\lean{Econ.sub_eq_integral_of_sandwich}
\leanok
If $f$ is monotone on $[a,b]$ and $f(x)(y-x) \le U(y) - U(x) \le f(y)(y-x)$ for $a \le x \le y \le b$, then
$U(b) - U(a) = \int_a^b f$.
\end{lemma}
\begin{proof}
\leanok
On a uniform partition with $N$ cells, telescoping and monotonicity place both $U(b)-U(a)$ and $\int_a^b f$ between the
lower and upper Riemann sums, whose gap is $(b-a)(f(b)-f(a))/N$. Let $N\to\infty$.
\end{proof}

\begin{definition}[Incentive compatible direct mechanism]
\label{def:ic-mech}
\lean{Econ.infoRent, Econ.IsICMech}
\leanok
Types $\theta \in [a,b]$; report $r$ gives allocation $q(r)$ for payment $p(r)$. IC: $\theta q(r) - p(r) \le
\theta q(\theta) - p(\theta)$ for all $\theta, r$. Rent $U(\theta) = \theta q(\theta) - p(\theta)$.
\end{definition}

\begin{theorem}[Myerson's lemma; MWG Prop.~23.D.2]
\label{thm:myerson-lemma}
\lean{Econ.myerson_ic_iff}
\leanok
\uses{def:ic-mech, lem:sandwich}
$(q,p)$ is IC on $[a,b]$ iff $q$ is nondecreasing and $U(\theta) = U(a) + \int_a^\theta q$ for all $\theta \in [a,b]$.
\end{theorem}
\begin{proof}
\leanok
IC for $x \le y$ both ways gives $q(x)(y-x) \le U(y) - U(x) \le q(y)(y-x)$; hence $q$ is monotone and the sandwich
lemma gives the envelope formula. Conversely $U(\theta) - U(r) = \int_r^\theta q \ge (\theta - r) q(r)$ by monotonicity,
which is IC.
\end{proof}

\begin{theorem}[Revenue equivalence, single-agent interim form; MWG Prop.~23.D.3]
\label{thm:revenue-equivalence}
\lean{Econ.revenue_equivalence, Econ.expected_revenue_eq}
\leanok
\uses{thm:myerson-lemma}
Two IC mechanisms with the same allocation rule and the same rent for the lowest type charge every type
the same payment, hence earn the same expected revenue under any type distribution on $[a,b]$.
\end{theorem}
\begin{proof}
\leanok
Both payments equal $\theta q(\theta) - U(a) - \int_a^\theta q$.
\end{proof}

\section{Optimal selling mechanism}
\label{sec:optimal-mechanism}

\begin{definition}[Virtual weights, discrete types]
\label{def:virtual-weight}
\lean{Econ.upperMass, Econ.virtualWeight, Econ.IsICDiscrete, Econ.IsIRDiscrete, Econ.rentFloor, Econ.postedQ, Econ.postedP}
\leanok
Types $\theta_0 \le \dots \le \theta_{n-1}$ with probabilities $f_j$. Virtual weight
$w_j = f_j\theta_j - (\theta_{j+1}-\theta_j)\sum_{k>j} f_k$ ($= f_j$ times the discrete virtual value).
IC and IR are imposed for all pairs of types. The posted price $\theta_m$ sells to types $j \ge m$ at price $\theta_m$.
\end{definition}

\begin{theorem}[Revenue identity and Myerson bound]
\label{thm:myerson-bound}
\lean{Econ.sum_mul_sum_range_lt, Econ.revenue_eq_virtual_sub, Econ.rentFloor_le, Econ.revenue_le_virtual}
\leanok
\uses{def:virtual-weight}
$\sum_k f_k p_k = \sum_j q_j w_j - \sum_k f_k\bigl(U_k - \sum_{i<k}(\theta_{i+1}-\theta_i)q_i\bigr)$; downward IC and IR
make the subtracted term nonnegative, so every IC and IR mechanism earns at most $\sum_j q_j w_j$.
\end{theorem}
\begin{proof}
\leanok
Exchange the order of summation in $\sum_k f_k \sum_{i<k}(\theta_{i+1}-\theta_i)q_i$; induct on $k$ for the rent floor.
\end{proof}

\begin{theorem}[Optimality of a posted price; Riley--Zeckhauser]
\label{thm:posted-price}
\lean{Econ.postedPrice_rent_eq, Econ.postedPrice_revenue, Econ.postedPrice_IC_IR, Econ.postedPrice_optimal, Econ.postedPrice_isOptimal}
\leanok
\uses{thm:myerson-bound}
A posted price attains the Myerson bound with equality and is IC and IR for nondecreasing types. If the
virtual weights are $\le 0$ below $m$ and $\ge 0$ from $m$ on, the posted price $\theta_m$ maximizes expected
revenue among all IC and IR mechanisms with $q \in [0,1]$.
\end{theorem}
\begin{proof}
\leanok
The posted-price rent equals the IC floor (telescoping), so the identity gives equality. For any feasible $q$,
$\sum_j q_j w_j \le \sum_{j \ge m} w_j$ term by term.
\end{proof}

\section{Bilateral trade}
\label{sec:bilateral-trade}

\begin{definition}[Bilateral trade mechanism]
\label{def:trade-mech}
\lean{Econ.TradeMechanism, Econ.TradeMechanism.IsEfficient, Econ.TradeMechanism.BuyerDSIC, Econ.TradeMechanism.SellerDSIC, Econ.TradeMechanism.IsBudgetBalanced}
\leanok
Buyer value $v$ and seller cost $c$ in $[a,b]$; a mechanism gives trade probability $x(v,c)$, buyer payment
$t_B$ and seller receipt $t_S$. Efficiency: $x = 1$ if $c < v$, $x = 0$ if $v < c$. Budget balance: $t_B = t_S$.
\end{definition}

\begin{theorem}[No efficient, strategy-proof, budget-balanced trade; Hurwicz, Green--Laffont]
\label{thm:bilateral-trade}
\lean{Econ.eq_of_squeeze, Econ.TradeMechanism.tB_eq_of_trade, Econ.TradeMechanism.tB_eq_of_noTrade, Econ.TradeMechanism.tS_eq_of_trade, Econ.TradeMechanism.tS_eq_of_noTrade, Econ.TradeMechanism.tB_sub_eq, Econ.TradeMechanism.tS_sub_eq, Econ.no_efficient_dsic_budgetBalanced_trade}
\leanok
\uses{def:trade-mech}
If $a < b$ no mechanism is efficient, DSIC for both traders, and budget balanced. (This is the
dominant-strategy counterpart of Myerson--Satterthwaite; the Bayesian version is not formalized.)
\end{theorem}
\begin{proof}
\leanok
At an interior seller report $c$, buyer IC pins the payment difference between trading and non-trading buyers
to exactly $c$; symmetrically at buyer report $v$ the seller's receipt difference is $v$. Take $x<y<z<w$ in $(a,b)$:
budget balance links $t_B(x,y) = t_S(x,y) = t_S(x,w) = t_B(x,w) = t_B(z,w) = t_S(z,w)$, and then
$y = t_B(z,y) - t_B(x,y) = t_S(z,y) - t_S(z,w) = z$, a contradiction.
\end{proof}
